\section{总结与延伸}

\subsection{全讲回顾}

这节课从 ``attention 是否足够'' 的问题出发，最终给出一整套 Transformer 的机制：
\begin{itemize}
    \item 用 self-attention 代替 recurrence，让交互距离降到 $O(1)$
    \item 用 positional encoding 把顺序信息重新注入模型
    \item 用 multi-head attention 捕捉不同类型的依赖
    \item 用 masked self-attention 保证 decoder 不偷看未来
    \item 用 cross-attention 让 decoder 访问 encoder 记忆
    \item 用 residual、LayerNorm 和 feedforward 让深层模型可训练
\end{itemize}

\begin{importantbox}{本讲的 takeaways}
\begin{enumerate}
    \item Transformers 不是只有 attention，而是 attention + 一整套配套设计
    \item self-attention 的最大优势是并行性和短路径依赖
    \item 位置编码是 Transformer 里不可或缺的一层
    \item 多头机制让模型可以同时学习多个关系视角
    \item 未来的改进主要集中在效率和位置建模上
\end{enumerate}
\end{importantbox}

\subsection{课程结语}

\begin{figure}[H]
\centering
\includegraphics[width=0.95\textwidth]{slides-images/slide-061.jpg}
\caption{Parting remarks：你现在已经理解了 Transformers，下一讲会进入 pre-training\protect\footnotemark}
\end{figure}
\footnotetext{Slides 第62页。}

讲义最后的提醒也很直接：继续准备 assignment 4，并关注后续的 project proposal。理解 Transformer 之后，下一步就是理解为什么 pre-training 能把模型性能再推高一个量级。

\begin{itemize}
    \item ``Good luck on assignment 4!''
    \item ``Remember to work on your project proposal!''
\end{itemize}

\subsection{延伸阅读}

\begin{itemize}
    \item Vaswani et al., Attention Is All You Need (2017): Transformer 的原始论文
    \item Shaw et al., Self-Attention with Relative Position Representations (2018): 相对位置编码
    \item Wang et al., Linformer / BigBird 等高效 attention 工作：降低 quadratic cost
    \item Vaswani 相关后续 lecture 和 CS224N 作业 4：Transformer 的实践训练
\end{itemize}

\end{document}
